\chapter{Signals}
\label{ch:signals}

This is the first chapter, so for now we can lean only on everyday
experience: a lamp is on or off, a door is open or closed. The question we
start with is simple: \emph{how do we turn ``the door is open'' into the zero
or one that a circuit can work with?} The answer comes from the notions of a
signal, logic levels, and the binary system, and along the way you will learn
to read timing diagrams.

\begin{infobox}[The Running Project ``Security Panel'']
Throughout Part I we will grow a single device: a security panel with a door
sensor, a window sensor, a night mode, an arming key, a code lock, a siren,
and an indicator. Each chapter adds a new block to it, and after every
chapter the panel still works. The first step is the smallest one: the door
sensor as a single bit and an ``alarm'' LED.
\end{infobox}

\begin{buildbox}
Open Logra, go to the \emph{Learn} tab, and complete the lessons ``A bit:
zero and one'' and ``The button and the LED on the board''. In the second
lesson, connect the button directly to the LED and press it. If the lamp did
not light up but went out instead, that is fine: by the end of the chapter you
will be able to explain why this happened and fix the circuit.
\end{buildbox}

\section{What Is a Signal}

A \term{signal} is a physical quantity that changes in time and carries
information. In electronics it is most often the \emph{voltage} on a wire. A
microphone turns sound into a changing voltage, a temperature sensor turns
temperature into a voltage, and a button turns a press into a voltage.

Signals come in two kinds.

\begin{description}
  \item[An analog signal] can take \emph{any} value within some range: 1.2~V,
    1.21~V, 1.215~V\ldots\ This is what sound looks like, or the signal from a
    microphone or a sensor.
  \item[A digital signal] takes only a few allowed values. In binary digital
    engineering there are two: \textbf{logic 0} (low voltage) and
    \textbf{logic 1} (high voltage).
\end{description}

\begin{figure}[H]
\centering
\begin{tikzpicture}
\begin{groupplot}[
  group style={group size=2 by 1, horizontal sep=18mm},
  width=0.5\textwidth, height=3.6cm, xmin=0, xmax=10,
  xlabel={Time}, axis lines*=left, xtick=\empty,
  label style={font=\sffamily\small}, tick label style={font=\sffamily\small},
  title style={font=\sffamily\small\bfseries},
]
\nextgroupplot[title={Analog signal}, ylabel={$U$, V}, ymin=0, ymax=3.5]
  \addplot[very thick, FpgaBlue, domain=0:10, samples=200] {1.65+1.2*sin(deg(x*0.9))+0.35*sin(deg(x*2.7))};
\nextgroupplot[title={Digital signal}, ymin=-0.2, ymax=1.4, ytick={0,1}, yticklabels={0,1}, ylabel={logic level}]
  \addplot[very thick, FpgaOrange, const plot] coordinates {(0,0) (1,1) (2.5,0) (3,1) (5,1) (5.5,0) (7,1) (8,0) (9,1) (10,1)};
\end{groupplot}
\end{tikzpicture}
\caption{An analog signal changes smoothly; a digital signal changes in jumps between two levels}
\end{figure}

\subsection{Why Digital}

Why is all modern computing digital? The main reason is \textbf{noise
immunity}. Any wire picks up noise: from neighboring wires, from motors, from
radio waves. In an analog signal noise distorts the information itself, and
it can no longer be separated from the useful signal. In a digital signal we
care about only one thing: whether the voltage is \emph{above or below a
threshold}. A small disturbance changes nothing.

\begin{figure}[H]
\centering
\begin{tikzpicture}
\begin{axis}[
  width=0.8\textwidth, height=5.4cm, xmin=0, xmax=10, ymin=-0.4, ymax=3.8,
  xlabel={Time}, ylabel={$U$, V}, xtick=\empty, ytick={0,0.8,2.0,3.3},
  yticklabels={0, 0.8, 2.0, 3.3},
  axis lines*=left, label style={font=\sffamily\small}, tick label style={font=\sffamily\small},
  clip=false,
]
  \fill[green!12] (axis cs:0,2.0) rectangle (axis cs:10,3.8);
  \fill[red!10] (axis cs:0,0.8) rectangle (axis cs:10,2.0);
  \fill[blue!8] (axis cs:0,-0.4) rectangle (axis cs:10,0.8);
  \addplot[thick, FpgaBlue, samples=400, domain=0:10]
    {(x>1 && x<4 || x>6 && x<8.5 ? 3.1 : 0.2) + 0.25*sin(deg(x*37)) + 0.15*sin(deg(x*91))};
  \node[font=\sffamily\small, anchor=west] at (axis cs:10.1,2.9) {``1''};
  \node[font=\sffamily\small, anchor=west, text=FpgaRed] at (axis cs:10.1,1.4) {forbidden};
  \node[font=\sffamily\small, anchor=west] at (axis cs:10.1,0.2) {``0''};
\end{axis}
\end{tikzpicture}
\caption{A signal with noise. As long as the voltage does not enter the forbidden zone,
the receiver tells 0 from 1 without error (the thresholds are for the 3.3~V LVCMOS standard)}
\label{fig:levels}
\end{figure}

\begin{infobox}[From Relay to Transistor]
The idea of ``on or off'' is older than electronics. The relay, a switch
controlled by an electromagnet, was used in telegraphy from the 1830s, and the
first computing machines were built from relays (the German Z3, 1941, with
about two thousand relays). A relay clicks in milliseconds, makes noise, and
wears out. ENIAC (1945) used about 17,500 vacuum tubes: faster, but tubes got
hot and burned out. In 1947 the transistor appeared at Bell Labs---the same
switch with no moving parts---and from 1958--1959 thousands of transistors
began to be made on a single chip. The Intel 4004 microprocessor (1971) had
about 2300 transistors, and modern chips have tens of billions. Only the
technology changed; the meaning stayed the same: each of these switches is in
one of two states, that is, it stores one bit.
\end{infobox}

\begin{diybox}[A Bit from a Sensor]
\textbf{Task.} Agree on what 0 and 1 mean for the door sensor: for example,
1 means ``the door is open'' and 0 means ``closed''. State it in one phrase:
such an agreement is called a level convention.\par
\textbf{In Logra.} The \emph{Learn} tab, the lesson ``A bit: zero and
one''.\par
\textbf{Expected result.} You confidently answer the question ``what does 1
on this wire mean'' without looking at the circuit.\par
\textbf{How to check.} Answer the questions of the lesson; the built-in check
will tell you whether you determined the value of the bit correctly.
\begin{itemize}[leftmargin=3.4em,labelwidth=3em,labelsep=0.4em,align=left]
  \item[\stars{1}] Repeat: complete the lesson to the end without hints.
  \item[\stars{2}] Modify: describe three more bits of your room (a window, a
    light switch, an arming key) and write down what 1 means for each.
  \item[\stars{3}] Project: a temperature sensor produces an analog voltage.
    Come up with a threshold (for example, 30\,\textdegree C) and explain how
    to get a digital ``too hot'' from the analog signal.
\end{itemize}
\end{diybox}

\section{Logic Levels}

Which voltage counts as zero and which as one is defined by a
\term{logic level standard}. The chip on the Tang Nano 20K learning board uses
the \textbf{LVCMOS33} standard (3.3~V supply). \lookahead{the board carries a
programmable chip (an FPGA), to which the second part of the book is
devoted.} The parameters of the standard:

\begin{table}[H]
\centering
\caption{LVCMOS 3.3 V levels}
\begin{tabular}{l l l}
\toprule
\textbf{Parameter} & \textbf{Value} & \textbf{Meaning} \\
\midrule
$U_\text{IL,max}$ & 0.8 V & anything below this, an input reads as zero \\
$U_\text{IH,min}$ & 2.0 V & anything above this, an input reads as one \\
$U_\text{OL,max}$ & 0.4 V & an output gives ``0'' at no more than this voltage \\
$U_\text{OH,min}$ & 2.4 V & an output gives ``1'' at no less than this voltage \\
\bottomrule
\end{tabular}
\end{table}

\begin{figure}[H]
\centering
\begin{tikzpicture}[yscale=1.3]
  \fill[blue!10] (0,0) rectangle (2.4,0.4);   \fill[green!15] (0,2.4) rectangle (2.4,3.3);
  \fill[gray!15] (0,0.4) rectangle (2.4,2.4);
  \fill[blue!10] (5,0) rectangle (7.4,0.8);   \fill[green!15] (5,2.0) rectangle (7.4,3.3);
  \fill[red!12] (5,0.8) rectangle (7.4,2.0);
  \draw[thick] (0,0) rectangle (2.4,3.3); \draw[thick] (5,0) rectangle (7.4,3.3);
  \node[font=\sffamily\small\bfseries] at (1.2,3.65) {Output};
  \node[font=\sffamily\small\bfseries] at (6.2,3.65) {Input};
  \node[font=\sffamily\small] at (1.2,2.85) {``1''}; \node[font=\sffamily\small] at (1.2,0.2) {``0''};
  \node[font=\sffamily\small] at (6.2,2.65) {``1''}; \node[font=\sffamily\small] at (6.2,0.4) {``0''};
  \node[font=\sffamily\fontsize{6}{7}\selectfont, text=FpgaRed] at (6.2,1.4) {undefined};
  \foreach \v/\l in {0.4/0.4, 2.4/2.4, 3.3/3.3} \node[font=\sffamily\scriptsize, anchor=east] at (0,\v) {\l~V};
  \foreach \v/\l in {0.8/0.8, 2.0/2.0} \node[font=\sffamily\scriptsize, anchor=west] at (7.4,\v) {\l~V};
  \draw[FpgaOrange, <->, thick] (3.7,2.0) -- (3.7,2.4) node[midway, right, font=\sffamily\scriptsize] {margin};
  \draw[FpgaOrange, <->, thick] (3.7,0.4) -- (3.7,0.8) node[midway, right, font=\sffamily\scriptsize] {margin};
  \draw[dashed, FpgaGray] (2.4,2.4) -- (5,2.4); \draw[dashed, FpgaGray] (2.4,2.0) -- (5,2.0);
  \draw[dashed, FpgaGray] (2.4,0.4) -- (5,0.4); \draw[dashed, FpgaGray] (2.4,0.8) -- (5,0.8);
\end{tikzpicture}
\caption{The output guarantees the levels ``with a margin'', so noise of 0.4~V
does not cause an error}
\end{figure}

\begin{infobox}[Active High and Active Low]
Usually the ``action'' happens at one: 1 means the LED is on, 1 means the
motor is running. This is \term{active high}. But often it is the other way
round: a reset signal \emph{fires at zero}. Such signals are called
\term{active-low} signals and are marked with an overbar
($\overline{\text{RESET}}$), a hash sign (RESET\#), or the suffix \code{\_n}
(\code{rst\_n}). On the Tang Nano 20K board the LEDs are exactly like this:
they light at zero.
\end{infobox}

\begin{diybox}[The Door Sensor and the ``Alarm'' LED]
\textbf{Task.} The door sensor gives 1 when the door is open. The board LED
lights at 0. Connect them so that the ``alarm'' LED lights when the door is
open. Hint: you will need an inverter (a NOT gate); we will look at how it
works in the next chapter.\par
\textbf{In Logra.} The \emph{Learn} tab, the lesson ``The button and the LED
on the board''; then the examples \code{button} (the button connected
directly to the LED) and \code{not} (with an inverter). Build your own
``input, inverter, output'' circuit on the \emph{Designer} tab.\par
\textbf{Expected result.} With door${}=1$ the LED is on, with
door${}=0$ it is off.\par
\textbf{How to check.} The \emph{Waveforms} tab: apply 0 to the input, then
1, and make sure that \code{led\_n} is always the opposite of the input.
\begin{itemize}[leftmargin=3.4em,labelwidth=3em,labelsep=0.4em,align=left]
  \item[\stars{1}] Repeat: in the \code{button} example set the button first
    to 0, then to 1, and explain why with the direct connection the board LED
    is on when the button is released.
  \item[\stars{2}] Modify: on the \emph{Pins} tab find the LED ports and
    determine from the name \code{led\_n} which level turns them on.
  \item[\stars{3}] Project. \emph{Stage of the ``Security Panel'' project: the
    door sensor as a bit.} Build the circuit ``door $\to$ NOT $\to$
    \code{led\_n}''. Acceptance criterion: on the waveform \code{led\_n} is an
    exact inversion of door over the whole interval, and the LED is on exactly
    when the door is open.
\end{itemize}
\end{diybox}

\section{The Binary Number System}

One digital wire carries one \term{bit}---0 or 1. To transmit a number larger
than one, we take several wires (bits) and use the \term{binary number
system}. Each digit position has a weight, a power of two:

\begin{figure}[H]
\centering
\begin{tikzpicture}
  \foreach \i/\b in {0/1,1/0,2/1,3/1} {
    \pgfmathtruncatemacro{\p}{3-\i}
    \pgfmathtruncatemacro{\w}{2^(3-\i)}
    \node[draw=FpgaBlue, fill=FpgaLight, thick, minimum size=11mm, font=\ttfamily\Large] (b\i) at (\i*1.6,0) {\b};
    \node[font=\sffamily\small, above=1mm of b\i] {$2^{\p}=\w$};
    \node[font=\sffamily\small, below=1mm of b\i] {$\b\cdot\w$};
  }
  \node[font=\sffamily\small, anchor=west] at (6.2,-0.85) {$= 8+0+2+1 = 11$};
  \node[font=\sffamily\small, anchor=west] at (6.2,0) {$1011_2 = 11_{10}$};
\end{tikzpicture}
\caption{Converting a binary number to decimal}
\end{figure}

With $n$ bits you can write $2^n$ different numbers: from $0$ to $2^n-1$.
Four bits give 16 combinations, eight (a \term{byte}) give 256. Long binary
numbers are hard to read, so they are written in the \term{hexadecimal}
system: every 4 bits are replaced by one digit from 0 to F.

\begin{table}[H]
\centering
\caption{The numbers from 0 to 15 in three number systems}
\small
\begin{tabular}{c c c @{\hspace{12mm}} c c c}
\toprule
\textbf{Dec.} & \textbf{Bin.} & \textbf{Hex.} & \textbf{Dec.} & \textbf{Bin.} & \textbf{Hex.} \\
\midrule
0 & \code{0000} & 0 & 8  & \code{1000} & 8 \\
1 & \code{0001} & 1 & 9  & \code{1001} & 9 \\
2 & \code{0010} & 2 & 10 & \code{1010} & A \\
3 & \code{0011} & 3 & 11 & \code{1011} & B \\
4 & \code{0100} & 4 & 12 & \code{1100} & C \\
5 & \code{0101} & 5 & 13 & \code{1101} & D \\
6 & \code{0110} & 6 & 14 & \code{1110} & E \\
7 & \code{0111} & 7 & 15 & \code{1111} & F \\
\bottomrule
\end{tabular}
\end{table}

\begin{figure}[H]
\centering
\begin{tikzpicture}
\begin{axis}[
  width=0.75\textwidth, height=5cm, ybar, bar width=6mm,
  xlabel={Number of bits $n$}, ylabel={Number of combinations $2^n$},
  xtick={1,...,8}, ymin=0, ymax=300,
  nodes near coords, nodes near coords style={font=\sffamily\scriptsize},
  label style={font=\sffamily\small}, tick label style={font=\sffamily\small},
  axis lines*=left, grid=major, grid style={FpgaGray!20},
]
  \addplot[fill=FpgaBlue!70, draw=FpgaBlue] coordinates {(1,2) (2,4) (3,8) (4,16) (5,32) (6,64) (7,128) (8,256)};
\end{axis}
\end{tikzpicture}
\caption{Each new bit doubles the number of possible values}
\end{figure}

What other number systems exist, how to write letters in binary, and how to
protect a number from errors during transmission we will cover in the next
chapter.

\begin{diybox}[A Number on LEDs]
\textbf{Task.} Output a four-bit number to four LEDs so that a one in a bit
position lights the LED. Remember that the board LEDs light at zero, so all
the bits will have to be inverted.\par
\textbf{In Logra.} The \emph{Learn} tab, the lesson ``A number on four
LEDs''.\par
\textbf{Expected result.} The number 5 (\code{0101}) lights the second and
fourth LEDs from the right, the number 0 lights nothing, and the number 15
lights all of them.\par
\textbf{How to check.} The built-in check of the lesson tries all 16
numbers; on the \emph{Waveforms} tab compare the output with the table of
numbers from 0 to 15 in this chapter.
\begin{itemize}[leftmargin=3.4em,labelwidth=3em,labelsep=0.4em,align=left]
  \item[\stars{1}] Repeat: complete the lesson and enter the numbers 5 and 11.
  \item[\stars{2}] Modify: which LEDs should be on for the number 13? Check
    your prediction on the waveform.
  \item[\stars{3}] Project: build from scratch in the \emph{Designer} a
    six-bit bus and output the number 37 to six LEDs (problem~2). Acceptance
    criterion: when 37 is entered, the LEDs of bits 0, 2, and 5 are on.
\end{itemize}
\end{diybox}

\section{Timing Diagrams}

To show how signals change in time, we draw \term{timing diagrams}: time runs
horizontally and the level vertically. A transition from 0 to 1 is called a
\term{rising edge} (or simply an edge), and a transition from 1 to 0 is a
\term{falling edge}.

\begin{figure}[H]
\centering
\begin{tikzpicture}[x=1cm, y=1cm]
  \draw[very thick, FpgaBlue] (0,0) -- (2,0) -- (2.2,1.2) -- (5,1.2) -- (5.2,0) -- (8,0);
  \draw[-{Stealth}, FpgaGray] (0,-1.0) -- (8.5,-1.0) node[right, font=\sffamily\small] {$t$};
  \node[font=\sffamily\small, anchor=east] at (0,0) {0};
  \node[font=\sffamily\small, anchor=east] at (0,1.2) {1};
  \draw[FpgaOrange, thick, -{Stealth}] (1.3,1.9) node[above, font=\sffamily\small] {rising edge ($0\to1$)} -- (2.05,0.8);
  \draw[FpgaOrange, thick, -{Stealth}] (6.3,1.9) node[above, font=\sffamily\small] {falling edge ($1\to0$)} -- (5.15,0.8);
  \draw[decorate, decoration={brace, amplitude=5pt, mirror}] (2.2,-0.1) -- (5,-0.1) node[midway, below=5pt, font=\sffamily\scriptsize] {pulse width};
\end{tikzpicture}
\caption{A pulse on a timing diagram}
\end{figure}

If several wires on a diagram are combined into a number, they are drawn as a
``bus'': two lines with the value written between them, and a cross marks the
moment when the value changes.

\begin{figure}[H]
\centering
\begin{tikztimingtable}[timing/slope=0.1, xscale=2.2, yscale=1.4, timing/font=\sffamily\small, timing/rowdist=3.2ex]
  a        & 2L 4H 2L 2H \\
  b        & 4L 4H 2L \\
  data[3:0] & 2D{0} 2D{5} 2D{A} 2D{F} 2D{3} \\
  z (disconnected) & 3Z 4H 3Z \\
\end{tikztimingtable}
\caption{Single signals, a bus, and a signal in the third (disconnected) state Z}
\end{figure}

The letter $Z$ on the last line means that the output is \emph{disconnected}
from the wire: it supplies neither a zero nor a one, and the wire may be taken
over by another source. How this works physically we will see later, in
Chapter~\ref{ch:electrical}.

\begin{diybox}[Reading a Waveform]
\textbf{Task.} Learn to read a timing diagram: find the rising and falling
edges, and measure the pulse width and the period.\par
\textbf{In Logra.} Open the example \code{blink} and the \emph{Waveforms}
tab.\par
\textbf{Expected result.} On the waveform of the LED signal you can point out
a rising edge, a falling edge, and two adjacent rising edges between which one
period passes.\par
\textbf{How to check.} Select two adjacent rising edges on the waveform and
compare the time difference with what you counted by eye.
\begin{itemize}[leftmargin=3.4em,labelwidth=3em,labelsep=0.4em,align=left]
  \item[\stars{1}] Repeat: find the period and the pulse width for the LED in
    the \code{blink} example.
  \item[\stars{2}] Modify: make the LED blink twice as fast (in the example
    code find the counter or constant that sets the speed) and compare the
    periods before and after.
  \item[\stars{3}] Project: complete the lesson ``A divide-by-2 frequency
    divider'' (we have not studied flip-flops yet; just follow the hints of
    the lesson). Acceptance criterion: the period of the output signal is
    exactly twice the period of the clock, and this is visible on the
    waveform (problem~6).
\end{itemize}
\end{diybox}

\section{The Clock Signal}

A special role in digital engineering is played by the \term{clock signal}
(\code{clk}): a sequence of pulses that arrive at strictly equal intervals of
time. It works like a conductor: on its edges all the memory elements of the
circuit update their values at the same moment.

\begin{figure}[H]
\centering
\begin{tikzpicture}[x=1cm, y=1cm]
  \foreach \i in {0,1,2,3} {
    \draw[very thick, FpgaBlue] (\i*3,0) -- (\i*3,1.2) -- (\i*3+1.5,1.2) -- (\i*3+1.5,0) -- (\i*3+3,0);
    \draw[FpgaOrange, -{Stealth}, thick] (\i*3,1.9) -- (\i*3,1.3);
  }
  \draw[-{Stealth}, FpgaGray] (-0.3,-0.4) -- (12.5,-0.4) node[right, font=\sffamily\small] {$t$};
  \draw[<->, thick] (3,-0.9) -- (6,-0.9) node[midway, below, font=\sffamily\small] {period $T$};
  \draw[<->, thick] (6,1.55) -- (7.5,1.55) node[midway, above, font=\sffamily\small] {$t_\text{w}$};
  \node[font=\sffamily\small, text=FpgaOrange, anchor=west] at (9.3,2.1) {rising edges};
\end{tikzpicture}
\caption{The clock signal: period $T$, pulse width $t_\text{w}$}
\end{figure}

The main characteristics of a clock signal:
\[
  f = \frac{1}{T}, \qquad D = \frac{t_\text{w}}{T}\cdot 100\,\%.
\]
Here $f$ is the \term{frequency} (how many periods per second, measured in
hertz), and $D$ is the \term{duty cycle} (usually 50\,\%). The Tang Nano 20K
board has a 27~MHz oscillator:
\[
  T = \frac{1}{27\,000\,000\ \text{Hz}} \approx 37 \cdot 10^{-9}\ \text{s} = 37\ \text{ns}.
\]
In one second 27 million clock cycles occur. During one blink of an eye
(about 0.1~s) there are almost three million.

\begin{table}[H]
\centering
\caption{Prefixes for frequency and time}
\begin{tabular}{l l l l}
\toprule
\textbf{Frequency} & \textbf{Period} & \textbf{Example} \\
\midrule
1 Hz  & 1 s   & an LED blinking \\
1 kHz = $10^3$ Hz & 1 ms = $10^{-3}$ s & an audible ``beep'' tone \\
1 MHz = $10^6$ Hz & 1 \textmu s = $10^{-6}$ s & simple microcontrollers \\
27 MHz & 37 ns & the Tang Nano 20K oscillator \\
1 GHz = $10^9$ Hz & 1 ns = $10^{-9}$ s & computer processors \\
\bottomrule
\end{tabular}
\end{table}

\begin{diybox}[Clock and Frequency]
\textbf{Task.} Relate the clock frequency to the period and to the time that
the events you need take.\par
\textbf{In Logra.} The \code{blink} example, the \emph{Waveforms} tab; the
formulas of this section.\par
\textbf{Expected result.} For any frequency in the prefix table you name the
period, and for a period you name the frequency.\par
\textbf{How to check.} Check your calculation against the prefix table; the
period on the waveform should match $1/f$.
\begin{itemize}[leftmargin=3.4em,labelwidth=3em,labelsep=0.4em,align=left]
  \item[\stars{1}] Repeat: find the period for the frequencies 1~kHz and 1~MHz.
  \item[\stars{2}] Modify: how many 27~MHz clock cycles pass in 1~ms? How many
    clock cycles are there in one period of a 1~Hz blink?
  \item[\stars{3}] Project: explain in your own words why an LED that
    switches every clock cycle at 27~MHz looks to us simply lit, and by how
    much the frequency must be divided to get a noticeable blink.
\end{itemize}
\end{diybox}

\begin{ideabox}
A digital signal has two levels: 0 and 1. Several bits form a binary number.
Changes of signals in time are shown with timing diagrams, and the clock
signal sets the rhythm at which the whole circuit works.
\end{ideabox}

\section{Common Mistakes}

\begin{warnbox}[Common Mistakes]
\begin{itemize}
  \item \textbf{Confusing ``active'' with ``lit''.} An active-low signal fires
    at zero. The board LED lights at 0, and a direct connection to a button
    that gives 1 when pressed turns the LED off when the button is pressed.
  \item \textbf{Reading a binary number left to right with the wrong weights.}
    The most significant bit is on the left, and the weight grows from right
    to left: $1011_2 = 8+0+2+1$, not $1+0+4+8$.
  \item \textbf{Thinking that $n$ bits give $n$ values.} They give $2^n$
    values, from $0$ to $2^n-1$; ten bits are 1024 numbers, not 10.
  \item \textbf{Trusting an ``in-between'' voltage.} A signal of 1.4~V at the
    input of 3.3~V logic is undefined: the circuit may read either 0 or 1.
\end{itemize}
\end{warnbox}

\section{Questions and Problems}
\label{sec:tasks-signals}

The answers are collected in Appendix~\ref{app:answers}.

\begin{taskbox}[Chapter Problems]
\begin{enumerate}
  \item \stars{1}~Convert the numbers $101101_2$ and $11111111_2$ to decimal.
  \item \stars{1}~Convert the numbers 37 and 200 to binary.
  \item \stars{1}~How many different values can be encoded with 10 bits? How many bits are needed to number 1000 objects?
  \item \stars{2}~For 3.3~V logic, a voltage up to 0.8~V is considered a low level and from 2.0~V a high level. What level does a 1.4~V signal have, and why are such signals undesirable?
  \item \stars{1}~What is the period of a 27~MHz clock signal? How many clock cycles pass in one second?
  \item \stars{2}~Draw a timing diagram of a clock signal and of a signal that changes on every rising edge (a divide-by-2 frequency divider). How many times slower is it?
  \item \stars{2}~\emph{Find the mistake.} A student connected a button (it gives 1 when pressed) directly to a board LED and complains that the LED is on when the button is released and goes out when it is pressed. Explain the cause and fix the circuit.
  \item \stars{2}~\emph{Find the mistake.} A student converted $1011_2$ to decimal like this: $1\cdot1+0\cdot2+1\cdot4+1\cdot8=13$. What is the mistake, and what is the correct answer?
  \item \stars{3}~\emph{Project.} Stage 1 of the ``Security Panel'' project: the door sensor gives 1 when the door is open, and the board LED lights at zero. Build in Logra a circuit in which the LED is on while the door is open. Acceptance criterion: on the \emph{Waveforms} tab the output \code{led\_n} is an exact inversion of the door input over the whole simulation interval.
  \item \stars{3}~\emph{Project.} Output the number 37 to six board LEDs (bits 5\ldots 0). Determine which LEDs must be on and what code must be applied to the \code{led\_n} outputs. Acceptance criterion: the circuit works for all numbers 0--63, and when 37 is entered the LEDs of bits 5, 2, and 0 are on.
\end{enumerate}
\end{taskbox}

\section{What's Next}

We have learned to record the state of a sensor with one bit and a number with
several bits, and to read waveforms. But a number in binary notation is not
the only way to represent information: the digits of a lock, letters, and the
position of a shaft are encoded in different ways, and transmitting a code
over a channel requires protection against errors. The next chapter,
Chapter~\ref{ch:codes}, is devoted to this. After that we will return to the
main question of the security panel: how to build the ``alarm'' rule from
several signals. For this we need logic gates (Chapter~\ref{ch:gates}).
